\subsection{LINEAR REPRESENTATIONS OF THE GROUP SL(2, k)}

Recall (Algebra, Chap. III, §8, no. 9) that we denote by \(\mathbf{SL}(2,k)\) the group of square matrices of order 2 with coefficients in \(k\) whose determinant is equal to 1. If \(x\in \mathfrak{sl}(2,k)\) is nilpotent, then \(x^{2} = 0\) (Algebra, Chap. VII, §5, Cor. 3 of Prop. 5) and \(e^x = 1 + x\in \mathbf{SL}(2,k)\) . If E is a finite dimensional vector space and \(\rho\) is a linear representation of \(\mathfrak{sl}(2,k)\) on E, then \(\rho (x)\) is nilpotent and so \(e^{\rho (x)}\) is defined (Chap. I, §6, no. 3).

DEFINITION 2. Let E be a finite dimensional vector space, and \(\rho\) (resp. \(\pi\) ) a linear representation of \(\mathfrak{sl}(2,k)\) (resp. \(\mathbf{SL}(2,k)\) ) on E. Then \(\rho\) and \(\pi\) are said to be compatible if, for every nilpotent element x of \(\mathfrak{sl}(2,k)\) , \(\pi(e^{x}) = e^{\rho(x)}\) .

In other words, \(\rho\) and \(\pi\) are compatible if, for every nilpotent element x of \(\mathfrak{sl}(2,k)\) , the restriction of \(\rho\) to kx is compatible with the restriction of \(\pi\) to the group \(1 + kx\) (Chap. VII, §3, no. 1).

If \(\rho\) and \(\pi\) are compatible, so are the dual representations, the mth tensor powers, and the mth symmetric powers of \(\rho\) and \(\pi\) , respectively (Chap. VII, §5, no. 4, Lemma 1 (i) and (ii)). Similarly for the representations induced by \(\rho\) and \(\pi\) on a vector subspace stable under \(\rho\) and \(\pi\) (loc. cit.).

In particular, the representation \(\rho_{m}\) of \(\mathfrak{sl}(2,k)\) on \(V(m)\) (no. 3) is compatible with the mth symmetric power \(\pi_{m}\) of the identity representation \(\pi_{1}\) of \(\mathbf{SL}(2,k)\) . Putting \(e_{n}^{(m)} = \binom{m}{n} u^{m-n} v^{n}\) as above, we have

\[
\pi_ {m} (s) e _ {n} ^ {(m)} = \binom {m} {n} (s u) ^ {m - n} (s v) ^ {n}\tag{3}
\]

for \(s \in \mathbf{SL}(2, k)\) and \(0 \leq n \leq m\) .

THEOREM 2. Let \(\rho\) be a linear representation of \(\mathfrak{sl}(2,k)\) on a finite dimensional vector space E.

\begin{enumerate}
\def\labelenumi{(\roman{enumi})}
\item
  There exists a unique linear representation \(\pi\) of \(\mathbf{SL}(2,k)\) on E that is compatible with \(\rho\) .
\item
  A vector subspace \(\mathrm{F}\) of \(\mathrm{E}\) is stable under \(\pi\) if and only if it is stable under \(\rho\) .
\item
  Let \(x \in \mathrm{E}\) . Then \(\pi(s)x = x\) for all \(s \in \mathbf{SL}(2,k)\) if and only if \(x\) is invariant under \(\rho\) (that is, \(\rho(a)x = 0\) for all \(a \in \mathfrak{sl}(2,k)\) ).
\end{enumerate}

The existence of \(\pi\) follows from the preceding and Th. 1. On the other hand, we know that the group \(\mathbf{SL}(2,k)\) is generated by the elements of the form

\[
e ^ {t X _ {+}} = \left( \begin{array}{c c} 1 & t \\ 0 & 1 \end{array} \right) \quad e ^ {- t X _ {-}} = \left( \begin{array}{c c} 1 & 0 \\ t & 1 \end{array} \right)
\]

where \(t \in k\) (Algebra, Chap. III, §8, no. 9, Prop. 17). This proves the uniqueness of \(\pi\) .

Assertions (ii) and (iii) follow from what we have said, together with Chap. VII, §3, no. 1, Lemma 1 (i). Q.E.D.

Every finite dimensional \(\mathfrak{sl}(2,k)\) -module therefore has a unique \(\mathbf{SL}(2,k)\) -module structure, which is said to be associated to its \(\mathfrak{sl}(2,k)\) -module structure.

Remark. When k is R or C or a complete non-discrete ultrametric field, \(\mathfrak{sl}(2,k)\) is the Lie algebra of \(\mathbf{SL}(2,k)\) . Let \(\rho\) and \(\pi\) be as in Th. 2. The homomorphism \(\pi\) is a homomorphism of Lie groups from \(\mathbf{SL}(2,k)\) to \(\mathbf{GL}(\mathrm{E})\) : this is clear when \(\mathrm{E} = \mathrm{V}(m)\) , and the general case follows, in view of Th. 1. By Chap. VII, §3, no. 1, \(\rho(X_{+}) = \mathrm{L}(\pi)(X_{+})\) , \(\rho(X_{-}) = \mathrm{L}(\pi)(X_{-})\) . Hence \(\rho = \mathrm{L}(\pi)\) (for a converse, see Exerc. 18).

PROPOSITION 5. Let E, F be finite dimensional \(\mathfrak{sl}(2,k)\) -modules, and let \(f \in \operatorname{Hom}_{k}(\operatorname{E},\operatorname{F})\) . The following conditions are equivalent:

\begin{enumerate}
\def\labelenumi{(\roman{enumi})}
\item
  \(f\) is a homomorphism of \(\mathfrak{sl}(2,k)\) -modules;
\item
  \(f\) is a homomorphism of \(\mathbf{SL}(2,k)\) -modules.
\end{enumerate}

Condition (i) means that \(f\) is an invariant element of the \(\mathfrak{sl}(2,k)\) -module \(\mathrm{Hom}_k(\mathrm{E},\mathrm{F})\) , and condition (ii) means that \(f\) is an invariant element of the \(\mathbf{SL}(2,k)\) -module \(\mathrm{Hom}_k(\mathrm{E},\mathrm{F})\) . Since these module structures are associated by Chap. VII, §5, no. 4, Lemma 1 (iii), the proposition follows from Th. 2 (iii).

DEFINITION 3. The adjoint representation of the group \(\mathbf{SL}(2,k)\) is the linear representation Ad of \(\mathbf{SL}(2,k)\) on \(\mathfrak{sl}(2,k)\) defined by

\[
\operatorname{Ad} (s). a = s a s ^ {- 1}
\]

for all \(a \in \mathfrak{sl}(2, k)\) and all \(s \in \mathbf{SL}(2, k)\) .

When k is R or C or a complete non-discrete ultrametric field, we recover Def. 7 of Chap. III, §3, no. 12 (cf.~loc. cit., Prop. 49).

By Chap. VII, §5, no. 4, Lemma 1 (i) and (ii), the adjoint representations of \(\mathfrak{sl}(2,k)\) and \(\mathbf{SL}(2,k)\) are compatible. By Chap. VII, §3, no. 1, Remark 2, \(\operatorname{Ad}(\mathbf{SL}(2,k)) = \operatorname{Aut}_e(\mathfrak{sl}(2,k))\) .

